\documentclass[11pt,a4paper]{article}

\usepackage{fontspec}
\setmainfont{DejaVu Serif}[
  BoldFont={DejaVu Serif Bold},
  ItalicFont={DejaVu Serif},
  ItalicFeatures={FakeSlant=0.18},
  BoldItalicFont={DejaVu Serif Bold},
  BoldItalicFeatures={FakeSlant=0.18}
]
\setsansfont{DejaVu Sans}
\usepackage{polyglossia}
\setmainlanguage{english}
\usepackage[a4paper,margin=20.5mm]{geometry}
\usepackage{microtype}
\usepackage{setspace}
\setstretch{1.008}
\usepackage{amsmath,amssymb,amsthm,mathtools,bm}
\usepackage{booktabs,array,tabularx}
\usepackage{enumitem}
\usepackage{xcolor}
\usepackage{tikz}
\usetikzlibrary{arrows.meta,positioning,calc}
\usepackage[unicode,hidelinks]{hyperref}
\usepackage[nameinlink,noabbrev]{cleveref}
\emergencystretch=3em
\allowdisplaybreaks

\hypersetup{
  pdftitle={Satellite Networks under Typed Frames and Temporal Observation Channels},
  pdfauthor={Stas, Independent Research Program},
  pdfsubject={Rigid-frame covariance, torus closures, line-of-sight temporal graphs, sampling, and relativistic clock protocols},
  pdfkeywords={satellite networks, observation geometry, SE(3), temporal graph, orbit closure, sampling alias, proper time, identifiability}
}

\definecolor{obsblue}{RGB}{37,83,138}
\definecolor{obsred}{RGB}{150,45,45}
\definecolor{obsgreen}{RGB}{30,115,78}
\definecolor{obsviolet}{RGB}{94,64,145}
\definecolor{lightblue}{RGB}{239,246,253}
\definecolor{lightred}{RGB}{253,242,242}
\definecolor{lightgreen}{RGB}{239,249,244}
\definecolor{lightviolet}{RGB}{246,242,252}

\theoremstyle{definition}
\newtheorem{definition}{Definition}[section]
\newtheorem{model}[definition]{Model}
\theoremstyle{plain}
\newtheorem{theorem}[definition]{Theorem}
\newtheorem{proposition}[definition]{Proposition}
\newtheorem{lemma}[definition]{Lemma}
\newtheorem{corollary}[definition]{Corollary}
\theoremstyle{remark}
\newtheorem{remark}[definition]{Remark}
\newtheorem{warning}[definition]{Boundary}

\newcommand{\R}{\mathbb{R}}
\newcommand{\Z}{\mathbb{Z}}
\newcommand{\Torus}{\mathbb{T}}
\newcommand{\SO}{\operatorname{SO}}
\newcommand{\SE}{\operatorname{SE}}
\newcommand{\rank}{\operatorname{rank}}
\newcommand{\diag}{\operatorname{diag}}
\newcommand{\clip}{\operatorname{clip}}
\newcommand{\dist}{\operatorname{dist}}
\newcommand{\cl}{\operatorname{cl}}
\newcommand{\spec}{\operatorname{spec}}
\newcommand{\norm}[1]{\left\lVert #1\right\rVert}
\newcommand{\abs}[1]{\left\lvert #1\right\rvert}
\newcommand{\Xcal}{\mathcal X}
\newcommand{\Ycal}{\mathcal Y}
\newcommand{\Gcal}{\mathcal G}
\newcommand{\Kcal}{\mathcal K}
\newcommand{\Qcal}{\mathcal Q}
\newcommand{\Scal}{\mathcal S}
\newcommand{\Acal}{\mathcal A}
\newcommand{\Ical}{\mathcal I}
\newcommand{\Ecal}{\mathcal E}
\newcolumntype{L}[1]{>{\raggedright\arraybackslash}p{#1}}
\newcolumntype{Y}{>{\raggedright\arraybackslash}X}

\title{\bfseries Satellite Networks under Typed Frames and Temporal Observation Channels\\
\large Rigid-frame covariance, torus closures, line-of-sight graphs, sampling, and clock protocols}
\author{Stas\\
\small Independent Research Program; ORCID: 0009-0000-2294-705X}
\date{Strict GO Core reference revision v1.2 --- 28 July 2026}

\begin{document}
\maketitle

\begin{center}
\fcolorbox{obsblue}{lightblue}{%
\parbox{0.92\linewidth}{\small
\textbf{Status.} This document supersedes the bilingual working note
v1.1 and is an application-layer reference migrated to GO Core v0.2.
It is not an orbit-determination method, an operational constellation
model, a replacement for TLE/SGP4, or a new result in astrodynamics.
Its contribution is a typed separation of hidden state, coordinate
frame, sensor channel, temporal graph, sampling protocol, closure,
and relativistic clock comparison.}}
\end{center}

\begin{abstract}
A satellite constellation is represented by one labeled dynamical
state but by many coordinate traces and still more sensor records.
The legacy note treated all of these as observer-frame drawings and
then called pairwise distances, coordinate spectra, orbit closures,
density statistics, and clock offsets ``invariants''.  This reference
states the transformation group and protocol for every such claim.
For a time-dependent rigid frame
\(\mathfrak f=(o,R)\in C^1(I,\SE(3))\), the labeled Euclidean distance
matrix is invariant at each common time, while component velocities,
coordinate spectra, and single-object orbit closures generally are
not.  A distance matrix determines a labeled configuration only up to
an element of \(E(3)\), so chirality is lost.  For a quasiperiodic
phase flow, the exact phase closure is the coset subtorus determined
by the integer resonance lattice; its image under a continuous
observation map is channel-dependent and is only equivariant under a
constant rigid isometry.  Finite sampling introduces exact frequency
aliases and cannot certify rational independence.  Spherical-body
line-of-sight and range gates define a temporal graph whose snapshot
and time-respecting connectivity are separated.  Proper time along a
fixed worldline segment is a spacetime scalar, whereas comparing
different clocks at equal coordinate time requires a declared
synchronization and gravity model.  These distinctions close the
last frame-law defect in the Geometry of Observation corpus.
\end{abstract}

\tableofcontents

\section{Scope, hidden state, and the observation chain}

Fix a coordinate-time interval \(I\subset\R\), a finite label set
\(V=\{1,\ldots,N\}\), and an Earth-centered inertial reference chart.
The minimal spatial state is the labeled tuple
\begin{equation}\label{eq:hidden-state}
X(t)=\bigl(c_E(t),(r_i(t),v_i(t))_{i\in V}\bigr),
\qquad v_i=\dot r_i ,
\end{equation}
where \(c_E\) is the center of the occulting body.  A set notation
\(\{r_i\}\) is insufficient: it discards labels and multiplicity, both
of which are required for tracking and temporal edges.

\begin{model}[Two-body benchmark]
\label{model:two-body}
The reproducibility controls use
\begin{equation}\label{eq:two-body}
\dot r_i=v_i,\qquad
\dot v_i=-\mu_E\frac{r_i-c_E}{\norm{r_i-c_E}^{3}},
\end{equation}
with a spherical exclusion radius \(R_E\).  This is an exact model
definition, not an ephemeris accuracy claim.
\end{model}

Classical elements \((a,e,i,\Omega,\omega,M_0)\) are a local
coordinate chart for \cref{model:two-body}; they are singular at
standard circular and equatorial limits and are not the hidden object.
Likewise, a two-line element set is a model-specific mean-element
record intended for a compatible SGP4-family propagator, not an
arbitrary Kepler state.  The former phrase ``TLE-like benchmark'' is
therefore removed.

\begin{definition}[Typed observation protocol]
A protocol is the composite
\begin{equation}\label{eq:chain}
\Xcal
\xrightarrow{\ \Phi_t\ }\Xcal_t
\xrightarrow{\ F_{\mathfrak f}\ }\Ycal_{\mathfrak f}
\xrightarrow{\ \Qcal\ }\mathcal D
\xrightarrow{\ \Scal_{\Delta t,T}\ }\mathcal D^M
\xrightarrow{\ \Acal\ }\Ical ,
\end{equation}
where \(\Phi_t\) is the declared dynamics,
\(F_{\mathfrak f}\) is an invertible coordinate change,
\(\Qcal\) is a possibly lossy sensor or graph channel,
\(\Scal_{\Delta t,T}\) samples over a finite horizon, and
\(\Acal\) is an estimator or diagnostic.  Information loss is
attributed to \(\Qcal\) or \(\Scal\), not to a frame change.
\end{definition}

\begin{center}
\begin{tikzpicture}[
  node distance=7mm and 8mm,
  box/.style={draw=obsblue,rounded corners,fill=lightblue,
    align=center,inner sep=4pt,font=\scriptsize},
  lossy/.style={draw=obsviolet,rounded corners,fill=lightviolet,
    align=center,inner sep=4pt,font=\scriptsize},
  arr/.style={-{Latex[length=2mm]},thick,draw=obsblue},
  larr/.style={-{Latex[length=2mm]},thick,draw=obsviolet},
  lab/.style={font=\tiny,fill=white,inner sep=1pt}
]
\node[box] (state) {labeled\\state};
\node[box,right=of state] (frame) {coframed\\coordinates};
\node[lossy,right=of frame] (channel) {range, bearing,\\LOS, graph};
\node[lossy,right=of channel] (sample) {sampled/noisy\\record};
\node[box,below=of channel] (inference) {protocol-scoped\\inference};
\draw[arr] (state) -- node[above,lab] {dynamics} (frame);
\draw[larr] (frame) -- node[above,lab] {channel} (channel);
\draw[larr] (channel) -- node[above,lab] {sampling} (sample);
\draw[arr] (sample) |- (inference);
\end{tikzpicture}
\end{center}

\section{Time-dependent rigid frames and exact spatial invariants}

Let \(o:I\to\R^3\) be a frame origin and let
\(R:I\to\SO(3)\) map frame components to inertial components.  The
coordinates of satellite \(i\) and of the Earth center are
\begin{equation}\label{eq:frame-map}
y_i=R^\mathsf T(r_i-o),\qquad
c_{\mathfrak f}=R^\mathsf T(c_E-o).
\end{equation}
Thus a frame is not merely an observer point.  Position vectors may
be composed across an Earth--Sun--barycentric tree only after all
origins, orientations, epochs, and time scales are declared.  In
homogeneous coordinates the legal composition is matrix
multiplication,
\begin{equation}\label{eq:frame-composition}
T_{A\leftarrow C}(t)
=T_{A\leftarrow B}(t)T_{B\leftarrow C}(t),
\end{equation}
not an untyped sum of vectors stored in different bases.

Let
\[
\Gcal_I=C^1(I,\SE(3))
\]
act diagonally on every node and on the central body through
\cref{eq:frame-map}.

\begin{theorem}[Distance-matrix invariance]
\label{thm:distance}
For every \(\mathfrak f\in\Gcal_I\), every common time \(t\), and
every \(i,j\in V\),
\begin{equation}\label{eq:distance-invariant}
D_{ij}^{\mathfrak f}(t)
=\norm{y_i(t)-y_j(t)}
=\norm{r_i(t)-r_j(t)}
=D_{ij}(t).
\end{equation}
Consequently, every deterministic function of the labeled distance
matrix is invariant under this declared action.
\end{theorem}

\begin{proof}
Translations cancel and \(R^\mathsf TR=I\):
\[
y_i-y_j=R^\mathsf T(r_i-r_j),\qquad
\norm{R^\mathsf T(r_i-r_j)}^2
=(r_i-r_j)^\mathsf T RR^\mathsf T(r_i-r_j).
\]
\end{proof}

The corresponding velocity law is different.  With
\(\Omega^\times=R^\mathsf T\dot R\in\mathfrak{so}(3)\),
\begin{equation}\label{eq:velocity-law}
\dot y_i=R^\mathsf T(v_i-\dot o)-\Omega^\times y_i .
\end{equation}
Therefore coordinate velocities, accelerations, and component spectra
are not invariants of \(\Gcal_I\).  Even the norm of a relative
velocity acquires the rotational term unless \(R\) is constant.

\begin{theorem}[What the distance matrix identifies]
\label{thm:edm}
Let \(D\) be the labeled Euclidean distance matrix of \(N\) points,
let \(J=I_N-\frac1N\bm1\bm1^\mathsf T\), and define
\begin{equation}\label{eq:gram}
B=-\frac12J(D\circ D)J .
\end{equation}
Then \(B\) is the centered Gram matrix.  If \(B\succeq0\) and
\(\rank B\le3\), its factorization reconstructs the configuration up
to \(E(3)\).  Pairwise distances do not distinguish translations,
rotations, or reflections.
\end{theorem}

\begin{proof}
Centering the points gives \(B=X_cX_c^\mathsf T\).  Any two
factorizations of the same positive semidefinite Gram matrix differ
by an orthogonal map on the realized span.  Restoring the centroid
adds an arbitrary translation.
\end{proof}

Thus the distance channel is complete for a labeled snapshot modulo
\(E(3)\), but not modulo \(\SE(3)\): chirality requires an oriented
volume or another parity-sensitive channel.

\section{Phase closure, image closure, and frame covariance}

For a quasiperiodic idealization, let
\begin{equation}\label{eq:phase-flow}
\theta(t)=\theta_0+\omega t\pmod{2\pi},
\qquad \theta\in\Torus^k ,
\end{equation}
and define the integer resonance lattice
\begin{equation}\label{eq:lattice}
L_\omega=\{m\in\Z^k:m\cdot\omega=0\}.
\end{equation}

\begin{theorem}[Exact phase closure]
\label{thm:torus}
The closure of the phase orbit is the coset subtorus
\begin{equation}\label{eq:phase-closure}
\Gamma_{\theta_0,\omega}
=\{\theta\in\Torus^k:
m\cdot(\theta-\theta_0)=0\pmod{2\pi}
\text{ for every }m\in L_\omega\},
\end{equation}
with
\[
\dim\Gamma_{\theta_0,\omega}
=k-\rank_{\Z}L_\omega .
\]
For every continuous \(F:\Torus^k\to\R^d\),
\begin{equation}\label{eq:image-closure}
\cl\{F(\theta(t)):t\in\R\}
=F(\Gamma_{\theta_0,\omega}).
\end{equation}
\end{theorem}

The rationally commensurate and rationally independent cases are only
the two endpoints of \cref{thm:torus}; intermediate resonance rank is
possible.  Compactness and continuity justify
\cref{eq:image-closure}.  A noncompact secular drift is not silently
treated as another torus angle.

\begin{proposition}[Closure is equivariant, not invariant]
\label{prop:closure}
For a constant \(g=(a,R)\in\SE(3)\) and
\(F_g(\theta)=a+RF(\theta)\),
\begin{equation}\label{eq:closure-equivariance}
K_{F_g}=gK_F .
\end{equation}
Equality \(K_{F_g}=K_F\) requires an additional symmetry of the set.
For a time-dependent \(g(t)\), even congruence may fail.
\end{proposition}

\begin{proof}
A homeomorphism commutes with closure, giving the first statement.
For the second, take \(F(\theta(t))\equiv0\) and
\(a(t)=(\cos\alpha t,\sin\alpha t,0)\).  The original closure is one
point and the transformed closure is a circle.
\end{proof}

The protocol-invariant object is therefore the phase-flow resonance
lattice, when the dynamical model and time parameter are fixed.  The
single-object spatial closure \(K_F\) is a channel output.

\section{Sampling, retarded observation, and non-identifiability}

At times \(t_m=t_0+m\Delta t\), \(m=0,\ldots,M-1\), the record is
finite.  Its support
\begin{equation}\label{eq:empirical-support}
\widehat K_{M,\Delta t}
=\{Q(X(t_m)):0\le m<M\}
\end{equation}
is not the infinite-time closure in \cref{eq:image-closure}.

\begin{proposition}[Exact sampling alias]
\label{prop:alias}
For any \(q\in\Z\) and \(t_m=m\Delta t\),
\begin{equation}\label{eq:alias}
e^{i(\omega+2\pi q/\Delta t)t_m}
=e^{i\omega t_m}.
\end{equation}
Hence a sampled coordinate spectrum does not identify an unrestricted
continuous frequency.
\end{proposition}

Finite noisy data also cannot decide whether a frequency ratio is
exactly rational: arbitrarily close rational and irrational vectors
produce arbitrarily close records on a fixed horizon.  Statements
about closure dimension therefore require either a declared exact
model or uncertainty-bounded inference.

The distinction between state and sensor time is equally important.
For a reception event \((t_r,o(t_r))\), a one-way ideal observation
uses an emission time \(t_{e,i}\) satisfying
\begin{equation}\label{eq:light-time}
t_r-t_{e,i}
=\frac{\norm{r_i(t_{e,i})-o(t_r)}}{c}.
\end{equation}
Different satellites are generally seen at different emission times.
Pairwise distances between apparent retarded positions are not the
simultaneous distance matrix in \cref{thm:distance}.

Bearing-only data
\[
b_i(t)=\frac{y_i(t)}{\norm{y_i(t)}}
\]
discard range.  Even without noise, one viewpoint therefore does not
recover \(D(t)\) unless additional baselines, ranging, or dynamical
hypotheses are supplied.

\section{Line-of-sight and temporal network geometry}

Let \(p_i=r_i-c_E\), \(d_{ij}=p_j-p_i\), and let the spherical
occulting radius including a declared margin be \(R_{\rm occ}>0\).
For \(i\ne j\), define
\begin{align}
u^\ast_{ij}
&=\clip\!\left(
-\frac{p_i\cdot d_{ij}}{\norm{d_{ij}}^2},0,1\right),\label{eq:clip}\\
h_{ij}
&=\norm{p_i+u^\ast_{ij}d_{ij}}.\label{eq:clearance}
\end{align}
The degenerate case \(p_i=p_j\) is handled separately.  A strict
spherical line of sight is declared when
\begin{equation}\label{eq:los}
\operatorname{LOS}_{ij}(t)
=\bm1\{h_{ij}(t)>R_{\rm occ}\}.
\end{equation}
Tangency belongs to the blocked class in this convention.

For link range \(R_{\rm link}>0\), the symmetric snapshot graph has
\begin{equation}\label{eq:adjacency}
A_{ij}(t)=
\bm1\{i\ne j\}
\bm1\{D_{ij}(t)\le R_{\rm link}\}
\operatorname{LOS}_{ij}(t).
\end{equation}
Because the satellites, the Earth center, and the segment are all
transformed together, \cref{eq:adjacency} is invariant under
\(\Gcal_I\).  A fixed spatial bin or an untransformed occulting center
would break that statement.

Under a vertex relabeling by a permutation matrix \(P\),
\[
A\longmapsto PAP^\mathsf T,\qquad
L_G=\diag(A\bm1)-A
\longmapsto PL_GP^\mathsf T.
\]
Connected-component counts and \(\spec(L_G)\) are label-invariant,
while node-specific routing outputs are equivariant under the same
relabeling.

\begin{definition}[Time-respecting journey]
Let an active edge \(e_k=(v_{k-1},v_k)\) be entered at \(t_k\) and
have a declared nonnegative transmission delay
\(\ell_{e_k}(t_k)\).  A journey satisfies
\begin{equation}\label{eq:journey}
t_{k+1}\ge t_k+\ell_{e_k}(t_k)
\end{equation}
for every consecutive edge.  Temporal reachability and earliest
arrival are defined by such journeys, not by connectivity of the
time-aggregated graph.
\end{definition}

The range-only graph
\(\bm1\{D_{ij}\le\rho\}\) is called a \emph{proximity graph}.
It is not a conjunction-risk calculation.  Operational collision
assessment additionally requires predicted relative states,
covariances, object sizes, maneuver hypotheses, and a time-window
protocol.

\section{Density and entropy require transported partitions}

Let \(K_\ell:\R^3\to[0,\infty)\) satisfy
\(\int_{\R^3}K_\ell(x)\,dx=1\), and let
\(w_i\ge0\), \(\sum_iw_i=1\).  The normalized occupancy density is
\begin{equation}\label{eq:density}
\rho_{\mathfrak f}(x,t)
=\sum_{i=1}^Nw_iK_\ell(x-y_i(t)).
\end{equation}
For a radial kernel and a rigid transform \(g=(a,R)\),
\(\rho_g(x,t)=\rho(R^\mathsf T(x-a),t)\).  This is covariance of a
field, not pointwise invariance at one fixed coordinate \(x\).

For a finite measurable partition
\(\Pi=\{B_1,\ldots,B_M\}\), define
\[
p_m=\int_{B_m}\rho(x,t)\,dx,\qquad
H_\Pi^{(2)}=-\sum_{p_m>0}p_m\log_2p_m .
\]
The entropy is preserved when \(\Pi\) is transported with the rigid
frame.  With fixed angular or Cartesian bins it is a
protocol-dependent diagnostic.  Kernel scale, partition, logarithm
base, horizon, and empty-total policy are mandatory fields.

\section{Proper time and the synchronization boundary}

Use metric signature \((-+++)\).  Along a fixed timelike worldline
segment \(\gamma_i\) with fixed endpoint events,
\begin{equation}\label{eq:proper-time}
\tau_i[\gamma_i]
=\frac1c\int_{\gamma_i}
\sqrt{-g_{\mu\nu}\,dx^\mu dx^\nu}.
\end{equation}
This scalar is invariant under coordinate changes.  In flat
spacetime, the corresponding event interval is invariant under the
Poincar\'e group.  Equal-time Euclidean distances are not Lorentz
invariants because a boost changes simultaneity.

In a declared stationary weak-field Earth chart with
\(\rho_i=\norm{r_i-c_E}\),
\(\Phi_i=-\mu_E/\rho_i\), and \(v\ll c\),
\begin{equation}\label{eq:weak-field}
\frac{d\tau_i}{dt}
=1+\frac{\Phi_i}{c^2}
-\frac{\norm{v_i}^2}{2c^2}
+O(c^{-4}).
\end{equation}
Relative to an ideal coordinate-stationary reference clock at
\(R_\star\),
\begin{equation}\label{eq:clock-offset}
\delta_i^{(t,\Phi)}
=\frac{d\tau_i}{dt}-\frac{d\tau_\star}{dt}
\simeq
\frac{\mu_E}{c^2}
\left(\frac1{R_\star}-\frac1{\rho_i}\right)
-\frac{\norm{v_i}^2}{2c^2}.
\end{equation}
For a circular two-body orbit,
\begin{equation}\label{eq:circular-offset}
\delta_{\rm circ}(r)
=\frac{\mu_E}{c^2}
\left(\frac1{R_\star}-\frac{3}{2r}\right).
\end{equation}

\begin{proposition}[A spatial locus does not determine proper time]
\label{prop:locus-clock}
An unparameterized spatial image \(K=\{r(t)\}\) does not determine
\(\tau\).
\end{proposition}

\begin{proof}
The same circle can be traversed with two different speed histories.
The spatial image is unchanged, while the kinetic term in
\cref{eq:weak-field} changes.  A time-stamped trajectory in a known
chart and a declared metric contains strictly more information than
its image.
\end{proof}

The comparison
\(\tau_i(T)-\tau_j(T)\) for endpoints selected by ``the same
coordinate time \(T\)'' is protocol-dependent: the time coordinate,
slicing, reference potential, and synchronization convention select
the endpoint events.  This does not contradict scalar invariance of
each fixed worldline integral.

\section{Benchmark controls and claim firewall}

The reproducibility bundle propagates a synthetic labeled network
under \cref{model:two-body}.  It checks time-dependent rigid-frame
distance invariance, Euclidean-distance-matrix reconstruction,
spherical segment clearance, graph invariance under co-transformed
geometry, Laplacian spectral invariance under relabeling, exact
sampling aliases, temporal earliest arrival, and weak-field clock
rates.  The numerical values are regression controls, not calibrated
constellation predictions.

\begin{table}[ht]
\centering
\small
\caption{Typed status of the principal legacy quantities.}
\begin{tabularx}{\linewidth}{L{30mm}L{34mm}Y}
\toprule
Object & Transformation or protocol & Correct status\\
\midrule
Distance matrix \(D(t)\) &
\(C^1(I,\SE(3))\), common time &
Invariant when every node is transformed diagonally.\\
Coordinate trace \(y_i(t)\) &
chosen frame \(\mathfrak f\) &
Covariant representation, not invariant.\\
Spatial closure \(K_F\) &
constant rigid isometry &
Equivariant up to congruence; generally changes for \(g(t)\).\\
Phase closure \(\Gamma_{\theta_0,\omega}\) &
fixed dynamics and time parameter &
Model-defined coset subtorus.\\
Coordinate spectrum &
frame plus sampling schedule &
Protocol-dependent and aliased.\\
LOS/link graph &
co-transformed body, fixed thresholds &
Snapshot invariant; temporal reachability needs ordered contacts.\\
Proper time \(\tau[\gamma]\) &
fixed worldline endpoints and metric &
Spacetime scalar.\\
Clock difference at \(t=T\) &
declared synchronization and potential &
Protocol-dependent comparison.\\
\bottomrule
\end{tabularx}
\end{table}

\begin{warning}[Claim boundary]
This reference proves elementary transformation, closure, sampling,
and graph statements inside declared models.  It does not provide
operational orbit determination, high-fidelity force models,
covariance propagation, collision probability, communication-link
budgets, clock steering, or a new theory of satellite networks.
\end{warning}

The minimum reproducible protocol records:
\begin{enumerate}[label=(\roman*),leftmargin=1.7em]
  \item labeled state variables, dynamics, constants, force model,
  epoch, reference frame, orientation convention, and time scale;
  \item observer origin and attitude, sensor type, light-time policy,
  occulting-body model, range/elevation masks, and noise model;
  \item sampling cadence, horizon, missing-data rule, quantizer or
  partition, estimator, tolerances, and uncertainty statement;
  \item for clocks, metric signature, potential convention,
  coordinate time, synchronization, endpoint events, and retained
  order in \(c^{-1}\).
\end{enumerate}

\section*{References}
\addcontentsline{toc}{section}{References}

\begin{enumerate}[label={[\arabic*]},leftmargin=2em,itemsep=2pt]
  \item G. Petit and B. Luzum (eds.),
  \emph{IERS Conventions (2010)}, IERS Technical Note 36,
  especially Chapters 5 and 10.
  \item CCSDS, \emph{Orbit Data Messages},
  CCSDS 502.0-B-3, Issue 3, April 2023.
  \item D. A. Vallado, P. Crawford, R. Hujsak, and T. S. Kelso,
  ``Revisiting Spacetrack Report \#3,''
  AIAA 2006-6753, \url{https://doi.org/10.2514/6.2006-6753}.
  \item GPS Directorate,
  \emph{Navstar GPS Space Segment/Navigation User Segment
  Interfaces}, IS-GPS-200N, 1 August 2022.
  \item N. Ashby, ``Relativity in the Global Positioning System,''
  \emph{Living Reviews in Relativity} 6, 1 (2003),
  \url{https://doi.org/10.12942/lrr-2003-1}.
  \item D. Kempe, J. Kleinberg, and A. Kumar,
  ``Connectivity and Inference Problems for Temporal Networks,''
  STOC 2000, 504--513,
  \url{https://doi.org/10.1145/335305.335364}.
  \item I. J. Schoenberg,
  ``Remarks to Maurice Frechet's Article,''
  \emph{Annals of Mathematics} 36 (1935), 724--732,
  \url{https://www.jstor.org/stable/1968654}.
  \item V. I. Arnold,
  \emph{Mathematical Methods of Classical Mechanics},
  2nd ed., Springer, 1989.
\end{enumerate}

\end{document}
